% Nouveau EGA V (cote 25): archive cote 25, site page 133 (PDF page 134).
% English translation of the checked French annotated master.
% AI-translated by GPT-6 Astra (OpenAI, ChatGPT Pro mode).
\scanpage{25/133}
\mspage{\hand{36}}

a k-vector subspace\var{\inserted[typed]{E}}{ E }{insertion}{ega25_25_133-v0001}of the space of \var{\deleted{\var{\inserted{pseudo}}{pseudo}{insertion}{ega25_25_133-v0003}}}{}{deletion}{ega25_25_133-v0002} \var{\deleted{rational}}{}{deletion}{ega25_25_133-v0004} \var{\inserted{meromorphic}}{meromorphic}{insertion}{ega25_25_133-v0005} \var{\deleted{\var{\inserted[typed]{strict}}{strict}{insertion}{ega25_25_133-v0007}}}{}{deletion}{ega25_25_133-v0006} functions \var{\xout{xxxxxxx}}{}{typewriter strike-out}{ega25_25_133-v0008}
on X, such that for every $\hand{\phi} \hand{\in} E-(0)$, $\hand{\phi}$ is \var{\inserted{regular}}{regular}{insertion}{ega25_25_133-v0009} \var{\deleted{invertible in the\var{\xout{xxxxxxx}}{}{typewriter strike-out}{ega25_25_133-v0011}}}{}{deletion}{ega25_25_133-v0010}
\var{\deleted{ring of rational functions}}{}{deletion}{ega25_25_133-v0012} \var{\deleted{\var{\inserted[typed]{that are strict}}{that are strict}{insertion}{ega25_25_133-v0014}}}{}{deletion}{ega25_25_133-v0013} \var{\xout{xxxxxx}}{}{typewriter strike-out}{ega25_25_133-v0015}, i.e. div($\hand{\phi}$) is defined \var{\deleted{as}}{}{deletion}{ega25_25_133-v0016}
\var{\deleted{a Cartier divisor,}}{}{deletion}{ega25_25_133-v0017}\hand{\unsure{,|)}} and $\hand{\triangle}$ is the set of div($\hand{\phi}$) for $\hand{\phi} \hand{\in} E-(0)$.
c) Let E, E' be two \var{\inserted[typed]{k-}}{k-}{insertion}{ega25_25_133-v0018}vector subspaces of the space of
rational functions \var{\inserted[typed]{that are strict}}{that are strict}{insertion}{ega25_25_133-v0019} \var{\xout{xxxxxxxx}}{}{typewriter strike-out}{ega25_25_133-v0020} on X, \var{\xout{xxxxx}}{}{typewriter strike-out}{ega25_25_133-v0021} satisfying the condition
considered in b); then the sets of divisors $\hand{\triangle}$, $\hand{\triangle}'$ that they
respectively define are equal iff there exists on X a \var{\inserted{pseudo-}}{pseudo-}{insertion}{ega25_25_133-v0022} function $\hand{\phi}$
\var{\deleted{strict rational}}{}{deletion}{ega25_25_133-v0023}e \var{\deleted{invertible}}{}{deletion}{ega25_25_133-v0024} \var{\inserted{that is regular}}{that is regular}{insertion}{ega25_25_133-v0025}, such that $E' = \hand{\phi} E$ . If $E\hand{\neq}(0)$\var{\inserted{(i.e. $\triangle \neq \emptyset$)}}{(i.e. $\triangle \neq \emptyset$)}{insertion}{ega25_25_133-v0026}, such
a $\hand{\phi}$ is determined modulo multiplication by an element c of $k^{\hand{*}}$.

The proof is an easy exercise, using \hand{4}.1. , and I
dispense with writing its proof, unless you ask me for it. It
\var{\margin{||}}{\placedmargin{||}}{marginal note placed}{ega25_25_133-v0027} seems to me moreover that \hand{4}.4. would advantageously come before \hand{4}.3.,
being technically distinctly more trivial. Observe moreover that when
X is geom integral, the condition on E stated in b) becomes \uline{vacuous}.
The restriction imposed at the beginning of b) comes from the fact that otherwise
the condition stated for b) might cease to hold after passage
to the algebraic closure of k (it is easy to give examples\var{\xout{xx}}{}{typewriter strike-out}{ega25_25_133-v0028},
in every characteristic, and even when k is separably closed in
c\var{\xout{\unsure{ar}}}{}{typewriter strike-out}{ega25_25_133-v0029} p \hand{$>$} 0). To do things properly, one should state b) without an additional condition
on k or X, but state the condition on E after passing to the
algebraic closure of k (and point out that when X is geom integral, \unsure{this}
condition becomes vacuous). \hand{-- By abuse of language, a}
\hand{\uline{set} $\triangle$ of divisors of the form \unsure{Ens|En}(D) will}
\hand{still sometimes be called a linear system of divisors}
\hand{on X ;}

\mspage{\hand{133}}
\tnote{25-133-1}{The isolated e after the canceled description of the function is a live source-letter residue. The source abbreviates geometrically as geom, and characteristic appears as c followed by the canceled doubtful ar and then p. The uncertainty translated as ``this'' has the French candidates ``oette'' and ``cette''. The manuscript's change from the capital delta on the preceding page to the triangle symbol here is retained.}
